\section{线段、凸组合与凸包}
\label{sec:09-Convex-Optimization/02-Convex-Sets/01}

假设你有两只股票，每只的持仓比例限制在 \(0\) 到 \(0.6\) 之间，且总持仓不超过 \(1\)。两个投资组合分别是

\[
x=(0.4,0.3),\qquad y=(0.2,0.5),
\]

它们各自都满足所有约束。

现在你想试试混合策略：把资金的 \(60\%\) 按 \(x\) 配置，\(40\%\) 按 \(y\) 配置。混合后的仓位是

\[
z=0.6\cdot x+0.4\cdot y=(0.32,0.38).
\]

检查一下：两只股票的仓位都在 \([0,0.6]\) 内，总和 \(0.7\le1\)——混合方案仍然可行。

这看起来理所当然，但换个情境试试。假设约束是``至多选择三只股票中的两只''。\(x\) 选了股票 A 和 B，\(y\) 选了 B 和 C，各自满足约束。但 \(\theta x+(1-\theta)y\) 在一个混合持仓里同时持有 A、B、C 三只——违反了``至多两只''的限制。混合操作在这里失效了。

两种情境的差别不在数字上，而在集合的形状上。前一个可行域里，任意两点的连线整段都在域内；后一个可行域有``凹陷''，连线会穿过不可行区域。

\subsection{凸集：线段永远不离开集合}

形式化上面的直觉，就得到了凸集的定义。集合 \(C\subseteq\R^n\) 是\textbf{凸集}，当且仅当

\[
x,y\in C,\ 0\le\theta\le1
\quad\Longrightarrow\quad
\theta x+(1-\theta)y\in C.
\]

定义只提到两点之间的线段，但它的力量可以延伸到任意有限个点。若 \(x_1,\ldots,x_k\in C\)，权重满足

\[
\theta_i\ge0,\qquad \sum_{i=1}^k\theta_i=1,
\]

那么

\[
\sum_{i=1}^k\theta_i x_i\in C.
\]

这样的点称为这些点的\textbf{凸组合}。证明用归纳法：先混合前 \(k-1\) 个点，再把结果与第 \(k\) 个点混合——每一步都只牵涉两点，凸性逐次保证结果留在集合内。

概率分布中的期望正是凸组合：\(p_i\) 是权重，\(x_i\) 是各可能取值。所以凸性与平均化、随机化天然绑在一起。

但代数上能取平均，不代表现实中一定能执行这个平均。若变量记录一周平均通勤时间与费用，五天里两天公交、三天骑车，就实现了两种方案 \(40\%\) 与 \(60\%\) 的平均——你的确可以在时间上混合。但今晚只能选一家餐厅，一顿饭不能按 \(40\%\) 川菜、\(60\%\) 日料来吃。可重复的行动允许按时间或按随机化混合；单次且不可分割的选择则不允许。

生产计划中同理。两种方案若可以按时间分配、且切换成本可忽略，分数安排确实可执行；若变量表示``是否建设一座工厂''或``选择哪一家唯一供应商''，凸组合只是原离散可行集的松弛：它能给下界、能帮助搜索，但要恢复真正可执行的方案，还需要舍入或组合决策。判断一个集合是否凸之前，先判断``混合''在变量的语义里到底代表什么。

\subsection{边界弯曲与凸性没有直接冲突}

欧氏单位球

\[
B=\{x:\norm{ x}_2\le1\}
\]

是凸集。取任意 \(x,y\in B\)，三角不等式给出

\[
\norm{\theta x+(1-\theta)y}_2
\le
\theta\norm{ x}_2+(1-\theta)\norm{ y}_2
\le1.
\]

所以弯曲的边界本身并不破坏凸性。开球、闭球都凸；凸性不要求集合是闭的、有界的，也不要求它包含原点。

反过来，圆周

\[
\{x:\norm{ x}_2=1\}
\]

不是凸集。取圆上两个相反的点，它们的中点是原点——原点不在圆周上。同理，两个互不相交圆盘的并集也不是凸集。

寻找两个集合内的点，使它们之间某个中间点落在集合外部——这是证明一个集合非凸最直接的方法。找一个反例就够了。有限采样只能帮你发现反例，不能用来证明一个无限集合是凸的。证明凸性需要覆盖任意 \(x,y\) 和任意 \(\theta\)，或者把集合表示为已知凸结构的组合。

\subsection{凸包：填平所有凹陷后的最小凸集}

给定任意集合 \(S\)——它可能到处都是洞和凹陷——它的\textbf{凸包}定义为所有有限凸组合的集合：

\[
\operatorname{conv}S = \left\{
\sum_{i=1}^k\theta_i x_i:
x_i\in S,\ \theta_i\ge0,\ \sum_i\theta_i=1
\right\}.
\]

这等价于``包含 \(S\) 的最小凸集''：任何包含 \(S\) 的凸集必须包含这些凸组合（否则不凸），而凸组合的凸组合仍然是凸组合，所以 \(\operatorname{conv}S\) 自身已经是凸集。

平面中三个不共线点的凸包是实心三角形，不是只有三条边的空心框。有限数据点的凸包给出了这些样本通过``混合''能够覆盖的全部区域——在几何上，它就是填平了所有凹陷之后的多面体。

概率单纯形

\[
\Delta^{n-1} = \{p\in\R^n:p_i\ge0,\ \mathbf1^{\trans}p=1\}
\]

是标准基向量 \(e_1,\ldots,e_n\) 的凸包。每个概率分布都可以看作``确定选择某个类别''的极端点之间的混合。线性目标在多面体的极端点上取得最优，几何原因正源于此：任何一个可行点都可以写成极端点的凸组合，而线性函数在凸组合上的值不会低于它在极端点上的最小值。

\subsection{哪些操作可靠地保持凸性}

许多实际集合不需要从定义出发重新证明。几条构造规则能覆盖绝大多数场景：

\textbf{任意族凸集的交仍然凸。} 若两点属于每个 \(C_\alpha\)，它们之间的线段也属于每个 \(C_\alpha\)，因而属于交集。优化问题中多个凸约束同时成立，正是取这些约束集合的交。

\textbf{仿射像保持凸性。} 若 \(C\) 凸，映射 \(x\mapsto Ax+b\) 的像

\[
\{Ax+b:x\in C\}
\]

仍凸，因为仿射映射与凸组合可交换：\(A(\theta x+(1-\theta)y)+b=\theta(Ax+b)+(1-\theta)(Ay+b)\)。

\textbf{仿射原像保持凸性。} 若 \(D\) 凸，则

\[
\{x:Ax+b\in D\}
\]

也是凸集。大量约束的凸性证明都可以通过这条规则一行完成——把约束写成某个已知凸集的原像。

\textbf{两个凸集的并一般不凸。} ``满足约束 A 或满足约束 B''往往引入离散选择，不能像``同时满足 A 与 B''那样保留凸性。``或''破坏凸性，``且''保留凸性——这一条在做建模取舍时反复出现。

投影也可看作一个线性像，因此把高维凸集投到较低维后仍然凸。这个事实很实用：消去一些辅助变量后，剩余变量的可行集合依然保持凸性，即使它的显式不等式并不容易写出来。

\subsection{一个容易误判的集合}

考虑正象限中的集合

\[
C=\{(x,y):xy\ge1,\ x>0,\ y>0\}.
\]

边界 \(y=1/x\) 向下弯，图形直觉可能会让人犹豫——它看起来``鼓''在外面。把约束改写为

\[
y\ge \frac1x,\qquad x>0.
\]

直接验证线段性质。任取 \((x_1,y_1),(x_2,y_2)\in C\) 和 \(\theta\in[0,1]\)，由 \(y_i\ge1/x_i\)，

\[
\theta y_1+(1-\theta)y_2
\ge\frac{\theta}{x_1}+\frac{1-\theta}{x_2}
\ge\frac{1}{\theta x_1+(1-\theta)x_2},
\]

即混合点仍满足 \(xy\ge1\)。最后一步是调和平均不超过算术平均：把
\(\bigl(\theta x_1+(1-\theta)x_2\bigr)\bigl(\theta/x_1+(1-\theta)/x_2\bigr)\) 展开，得到 \(1+\theta(1-\theta)(x_1-x_2)^2/(x_1x_2)\ge1\)。

用下一单元的语言，这正是``凸函数 \(1/x\) 的上境图是凸集''的一次手工验证；上境图判据届时会把这类论证压缩成一行。这个例子想说明的是：``看起来鼓或凹''不够可靠。凸集合的判断要回到线段、构造规则或函数上境图——视觉直觉容易出错，代数计算才不会骗你。
